\section{Sovereign Business Brain: Remaining Release Gates}

This section is the closure register for work that remains after the internally buildable
Sovereign Business Brain foundations completed on 5 September 2026.  An item remains here when it
requires a publisher identity, a customer-controlled environment, an accredited or independent
reviewer, a production provider account, representative hardware, or real customer evidence.  No
developer artefact is described as a qualified consumer release merely because its implementation
and local tests pass.

\subsection{Native distribution and clean-machine qualification}

The unsigned macOS Apple-Silicon package, Windows x64 developer archive and Linux x64 systemd
archive are complete.  The remaining release work is:
\begin{itemize}
  \item Apple Developer ID signing, notarisation and stapling, followed by clean Intel and Apple
        Silicon installation, upgrade, rollback and removal tests;
  \item Windows Authenticode signing and equivalent clean-machine Service Control Manager tests;
  \item signed native Linux packages and qualification across the declared systemd distributions;
  \item a published cross-platform matrix covering installation, recovery, update, rollback,
        backup, export, removal and network-disconnected operation.
\end{itemize}
These claims require the relevant publisher identities and representative target systems.

\subsection{Private compute and capability qualification}

The model router, compute contracts, capacity planner and signed capability-package format are
implemented.  Production qualification still requires one authorised customer-owned
OpenAI-\allowbreak compatible model endpoint and one measured NVIDIA NIM/NeMo deployment. Each must prove
context minimisation, endpoint identity, residency, time and cost ceilings, failure fallback,
result validation and provider replacement without loss of the AION brain.

Arbitrary third-party capability code also needs a production OS or container sandbox.  The six
completed core capability packs must be exercised against authorised real adapters and customer
outcomes before they can be advertised as production-qualified.  Their current signatures are
development identities and their manifests state that limitation.

\subsection{Commercial completion}

The entitlement, department-trial, cancellation, verified-action metering, capacity-control,
private-value-ledger, Boardroom and signed-phone surfaces, and opaque payment-boundary foundations
are complete.  Two commercial gates remain:
\begin{enumerate}
  \item connect a certified payment provider using only opaque vault references, exact visible
        terms and signed, idempotent provider events;
  \item validate packaging, conversion, limits and willingness to pay with real customers.
\end{enumerate}
Pilot Mobile derives the business from a signed, active canonical organisation membership and
never accepts an arbitrary tenant identifier or infers a business from a shared screen.

\subsection{Independent and operational assurance}

The remaining assurance cannot be self-certified by the implementation team.  It comprises an
independent penetration and software-supply-chain review; qualified privacy, employment, sector
and AI-regulation assessment for each launch region; operational service-level, disaster-recovery,
incident-response and vulnerability-disclosure exercises; and representative pilots with both a
small business and a multi-unit organisation.

\subsection{Release order and stop conditions}

The required order is certified payment-provider integration, platform signing and clean-machine
qualification, live private-compute and capability-pack
qualification, and finally independent assurance and field trials.  Any failed signature,
membership, residency, disclosure, evidence, payment-event or rollback check fails closed.  The
customer's brain, encrypted memory, Business Map, local models, backup and export remain available
when a paid entitlement expires; only the paid operational capability is reduced.
